% \iffalse meta-comment
%
% hit-footnote.dtx 是各个类共用的脚注层
%
% \fi
%
% \section{共享脚注层}
%
% \changes{v3.2a}{0000/00/00}{脚注层抽出来给各个类共用，art 与 artplus 的脚注
%   由此改成带圈编号、小五号宋体}
% 我工规范对脚注的要求各个类是一样的：编号带圈、按页重编、宋体小五、悬挂 1.5 字符。
% 原来这套只写在 \file{hithesisbook.cls} 里，art 与 artplus 用的还是 \LaTeX{} 默认的
% 上标阿拉伯数字。
%
% 分成两截。\cs{hit@setup@footnote@pkg} 只管加载 \pkg{footmisc}，位置要留在依赖模块
% 里不能动；\cs{hit@setup@footnote} 管余下的样式，得排在 \pkg{footmisc} 与字号都就位
% 之后调用，因为它要接管 \cs{@makefntext} 并存一份当时的 \cs{footnotesize}。
%
%    \begin{macrocode}
%<*shared-footnote>
%    \end{macrocode}
%
% \begin{macro}{\hit@setup@footnote@pkg}
% 禁止 \LaTeX 自动调整多余的页面底部空白，并保持脚注仍然在底部。脚注按页编号。
%    \begin{macrocode}
\cs_new_protected:Npn \hit@setup@footnote@pkg
  {
    \bool_if:NTF \g__hit_raggedbottom_bool
      {
        \RequirePackage [ bottom , perpage , hang ] { footmisc }
        \raggedbottom
      }
      { \RequirePackage [ perpage , hang ] { footmisc } }
  }
%    \end{macrocode}
% \end{macro}
%
% \begin{macro}{\hit@setup@footnote}
%    \begin{macrocode}
\cs_new_protected:Npn \hit@setup@footnote
  {
%    \end{macrocode}
%
% \changes{v3.2a}{0000/00/00}{带圈脚注号改用 \pkg{circledsteps} 的 \cs{Circled}}
% 带圈脚注号原来用 \cs{textcircled}，它是拿字体里的组合用圈套上数字，圈偏、数字
% 也不居中。改用 \pkg{circledsteps} 的 \cs{Circled}，圈是 Ti\emph{k}Z 画的，位置
% 正、两位数还会自动加宽。
%    \begin{macrocode}
\RequirePackage{circledsteps}
%    \end{macrocode}
%
% \changes{v3.2a}{0000/00/00}{带圈数字支持多位数，圈保持正圆}
% \cs{Circled} 里字数一多圈就拉成椭圆。量一下内容比一位数宽多少，按这个比例横向
% 压回去，圈就始终是正圆。原来那个模板用的是固定的 \texttt{0.45}，那是因为它的
% 西文数字取自 SimSun、是全宽的；这里西文是比例宽度，固定常数不适用，改成算出来。
%
% \cs{quan} 一到两位、\cs{Quan} 任意位数、\cs{quansplit} 按位数分派——三个都保留，
% 但压缩比例自动算，所以实际走的是同一条路。
%    \begin{macrocode}
\box_new:N \l__hit_quan_box
\box_new:N \l__hit_quanmark_box
\int_new:N \l__hit_quan_slot_int
\dim_new:N \l__hit_quandia_dim
\fp_new:N \l__hit_quank_fp
\fp_new:N \l__hit_quanx_fp
%    \end{macrocode}
%
% 圈里数字的横向压缩比例，与圈线的粗细（相对 em）。两个常数都是拿
% \option{fontset}\texttt{=windows}（SimSun）的自带带圈字符量出来的：数字压到
% 0.45 倍看着与字形版一致；字形的圈线粗细约为 0.034 em，而
% \cs{linethickness} 的默认值是绝对长度，字号一放大就明显偏细，所以改成跟着 em 走。
%    \begin{macrocode}
\fp_const:Nn \c__hit_quan_fill_fp { 0.5 }
\dim_const:Nn \c__hit_quan_rule_dim { 0.045em }
%    \end{macrocode}
%
% 把当前 CJK 字体设成当前字体。\cs{songti} 只是选 CJK 家族，西文数字仍走西文字体；
% 圈号里的数字要跟正文宋体一致，得真的切过去。
%    \begin{macrocode}
\cs_new_protected:Npn \__hit_cjkfont:
  { \usefont { \f@encoding } { \CJK@family } { m } { n } }
%    \end{macrocode}
%
% \texttt{n} 对应的带圈数字码位：1--20 在 \texttt{U+2460}，21--35 在
% \texttt{U+3251}，36--50 在 \texttt{U+32B1}，超出范围返回 0。
%    \begin{macrocode}
\cs_new:Npn \__hit_quan_slot:n ##1
  {
    \int_compare:nNnTF {##1} < { 1 } { 0 }
      {
        \int_compare:nNnTF {##1} < { 21 } { \int_eval:n { ##1 - 1 + "2460 } }
          {
            \int_compare:nNnTF {##1} < { 36 } { \int_eval:n { ##1 - 21 + "3251 } }
              {
                \int_compare:nNnTF {##1} < { 51 }
                  { \int_eval:n { ##1 - 36 + "32B1 } } { 0 }
              }
          }
      }
  }
%    \end{macrocode}
%
% 当前宋体有现成的带圈字符就直接用，字形比画出来的规整；没有才画。
%
% 画的几何按字体自带的 \texttt{U+2460} 定：量出它的直径，数字缩到直径的一半（字形
% 版里数字大致就占这个比例），剩下的补成 \pkg{circledsteps} 的内边距，圈就与字形版
% 一样大、数字也一样大。位数多了只在横向压，圈保持正圆。圈线粗细跟 em 走——圈线粗细
% 取 0.045~em，是拿 SimSun 的字形量出来的；\cs{linethickness} 的默认值是绝对长度，
% 字号一放大就偏细，所以跟着 em 走。
% \changes{v3.2a}{0000/00/00}{\cs{quan} 重做：先取当前宋体的带圈字符，超出字体
%   可用范围才画圈；画出来的几何按字体自带的 \texttt{U+2460} 定，圈线粗细跟 em 走。
%   原先的 \cs{Quan} 一并去掉，两种写法合成一个}
%    \begin{macrocode}
\NewDocumentCommand \quan { m }
  {
    \group_begin:
      \__hit_cjkfont:
      \int_set:Nn \l__hit_quan_slot_int { \__hit_quan_slot:n {##1} }
      \bool_lazy_and:nnTF
        { \int_compare_p:nNn \l__hit_quan_slot_int > { 0 } }
        { \int_compare_p:nNn { \XeTeXcharglyph \l__hit_quan_slot_int } > { 0 } }
        { \tex_char:D \l__hit_quan_slot_int }
        { \__hit_quan_draw:n {##1} }
    \group_end:
  }
\cs_new_protected:Npn \__hit_quan_draw:n ##1
  {
    \hbox_set:Nn \l__hit_quanmark_box { \tex_char:D "2460 }
    \hbox_set:Nn \l__hit_quan_box { ##1 }
    \dim_set:Nn \l__hit_quandia_dim
      { \box_ht:N \l__hit_quanmark_box + \box_dp:N \l__hit_quanmark_box }
    \dim_compare:nNnT { \l__hit_quandia_dim } = { 0pt }
      { \dim_set:Nn \l__hit_quandia_dim { 1em } }
%    \end{macrocode}
%
% 数字缩到直径的一半；横向如果还塞不进内圈，再按比例压窄。
%    \begin{macrocode}
    \fp_set:Nn \l__hit_quank_fp
      {
        \c__hit_quan_fill_fp * \dim_to_fp:n { \l__hit_quandia_dim }
        / \dim_to_fp:n { \box_ht:N \l__hit_quan_box + \box_dp:N \l__hit_quan_box }
      }
    \fp_set:Nn \l__hit_quanx_fp
      {
        \fp_min:nn { 1 }
          {
            \c__hit_quan_fill_fp * \dim_to_fp:n { \l__hit_quandia_dim }
            / ( \l__hit_quank_fp * \dim_to_fp:n { \box_wd:N \l__hit_quan_box } )
          }
      }
    \group_begin:
      \linethickness { \c__hit_quan_rule_dim }
      \pgfkeys
        {
          /csteps/inner~ysep =
            \fp_to_dim:n
              { \dim_to_fp:n { \l__hit_quandia_dim } * ( 1 - \c__hit_quan_fill_fp ) } ,
          /csteps/inner~xsep =
            \fp_to_dim:n
              { \dim_to_fp:n { \l__hit_quandia_dim } * ( 1 - \c__hit_quan_fill_fp ) } ,
        }
      \Circled
        {
          \scalebox { \fp_eval:n { \l__hit_quank_fp * \l__hit_quanx_fp } }
            [ \fp_use:N \l__hit_quank_fp ] { \box_use:N \l__hit_quan_box }
        }
    \group_end:
  }
\cs_new_protected:Npn \hit@textcircled ##1 { \quan { \arabic {##1} } }
\cs_set:Npn \thefootnote   { \hit@textcircled { footnote } }
\cs_set:Npn \thempfootnote { \hit@textcircled { mpfootnote } }
%    \end{macrocode}
%
% 定义脚注分割线，字号（宋体小五），以及悬挂缩进（1.5字符）。
%    \begin{macrocode}
\cs_set:Npn \footnoterule
  { \vskip -3\p@ \hrule \@width 0.3\textwidth \@height 0.4\p@ \vskip 2.6\p@ }
\cs_set_eq:NN \hit@footnotesize \footnotesize
\cs_set:Npn \footnotesize { \hit@footnotesize \xiaowu [1.5] }
\footnotemargin1.5em\relax
%    \end{macrocode}
%
% \cs{@makefnmark} 默认是上标样式，而在脚注部分要求为正文大小。
% \changes{v3.2a}{0000/00/00}{\cs{@makefntext} 改成包一层，不再用 \pkg{etoolbox}
%   前后插码。\cs{pretocmd} 会把宏体 detokenize 后按当前 catcode 重扫，在 expl3
%   里会吃掉 \cs{@makefntext} 体内的字面空格，脚注整块上移 4.2bp}
% 原来是往 \cs{@makefntext} 的宏体前后各插一句 \cs{let}，效果等于在原宏体跑之前
% 换掉 \cs{@makefnmark}、跑完再换回来。包一层做的是同一件事，位置也一样。
%    \begin{macrocode}
\cs_set_eq:NN \hit@makefnmark \@makefnmark
\cs_new:Npn \hit@@makefnmark { \hbox { { \normalfont \@thefnmark } } }
\cs_set_eq:NN \hit@orig@makefntext \@makefntext
\cs_set:Npn \@makefntext ##1
  {
    \cs_set_eq:NN \@makefnmark \hit@@makefnmark
    \hit@orig@makefntext {##1}
    \cs_set_eq:NN \@makefnmark \hit@makefnmark
  }
  }
%    \end{macrocode}
% \end{macro}
%
%    \begin{macrocode}
%</shared-footnote>
%    \end{macrocode}
%
% \Finale
